\documentclass[
    language=english,
    doctype=article,
    institution=none,
    titlestyle=book,
]{../../omnilatex}

\RequirePackage{config/document-settings}
\addbibresource{../../bib/bibliography.bib}

\RequirePackage{amsmath,amssymb,amsthm}
\RequirePackage{mathtools}
\RequirePackage{cancel}
\RequirePackage{enumitem}
\RequirePackage{tikz-cd}
\RequirePackage{booktabs}
\RequirePackage{subcaption}
\RequirePackage{float}

\DeclareMathOperator{\tr}{tr}
\DeclareMathOperator{\rank}{rank}
\DeclareMathOperator{\diag}{diag}
\DeclareMathOperator{\sgn}{sgn}
\DeclareMathOperator{\adj}{adj}
\DeclareMathOperator{\Hom}{Hom}
\DeclareMathOperator{\End}{End}
\DeclareMathOperator{\colim}{colim}
\DeclareMathOperator{\Var}{Var}
\DeclareMathOperator{\Cov}{Cov}
\DeclareMathOperator{\supp}{supp}
\DeclareMathOperator{\Res}{Res}
\DeclareMathOperator{\erf}{erf}
\DeclareMathOperator{\Lap}{\mathcal{L}}
\DeclareMathOperator{\iLap}{\mathcal{L}^{-1}}
\DeclareMathOperator{\Gal}{Gal}
\DeclareMathOperator{\lcm}{lcm}
\DeclareMathOperator{\GL}{GL}
\DeclareMathOperator{\SL}{SL}
\DeclareMathOperator{\SO}{SO}
\DeclareMathOperator{\im}{im}
\DeclareMathOperator{\ker}{ker}
\DeclareMathOperator{\Id}{Id}

\theoremstyle{plain}
\newtheorem{theorem}{Theorem}[section]
\newtheorem{lemma}[theorem]{Lemma}
\newtheorem{proposition}[theorem]{Proposition}
\newtheorem{corollary}[theorem]{Corollary}

\theoremstyle{definition}
\newtheorem{definition}{Definition}[section]
\newtheorem{example}{Example}[section]

\theoremstyle{remark}
\newtheorem{remark}{Remark}[section]

\begin{document}

\maketitle

\begin{abstract}
This gallery demonstrates 20 advanced mathematical typesetting techniques for
scientific publications, covering commutative diagrams, tensor notation, category
theory, proof environments, aligned equations, systems of equations, matrices,
operator definitions, dimensional analysis, chemical equations, logic, set theory,
probability, statistics, optimization, control theory, Fourier analysis,
differential geometry, algebraic structures, and number theory.
\end{abstract}

\tableofcontents

% =====================================================
\section{Commutative Diagrams}
\label{sec:commutative-diagrams}
% =====================================================

\subsection{Exact Sequences}

A short exact sequence of groups:
\begin{equation}\label{eq:short-exact}
    0 \longrightarrow A \xrightarrow{f} B \xrightarrow{g} C \longrightarrow 0
\end{equation}

The long exact sequence in homology:
\[
    \cdots \to H_n(A) \xrightarrow{f_*} H_n(B) \xrightarrow{g_*} H_n(C)
    \xrightarrow{\partial} H_{n-1}(A) \to \cdots
\]

\subsection{Pushouts and Pullbacks}

The pushout of $f\colon A \to B$ and $g\colon A \to C$:
\[
    \begin{tikzcd}
        A \arrow[r, "f"] \arrow[d, "g"'] & B \arrow[d, "h"] \\
        C \arrow[r, "k"'] & P \arrow[ul, phantom, "\ulcorner" near end]
    \end{tikzcd}
\]

The pullback of $f\colon B \to A$ and $g\colon C \to A$:
\[
    \begin{tikzcd}
        Q \arrow[r, "p"] \arrow[d, "q"'] \arrow[dr, phantom, "\lrcorner" near end]
            & C \arrow[d, "g"] \\
        B \arrow[r, "f"'] & A
    \end{tikzcd}
\]

\subsection{Natural Transformations}

A natural transformation $\eta\colon F \Rightarrow G$ satisfies naturality for every
$f\colon X \to Y$ in $\mathcal{C}$:
\[
    \begin{tikzcd}
        F(X) \arrow[r, "F(f)"] \arrow[d, "\eta_X"'] & F(Y) \arrow[d, "\eta_Y"] \\
        G(X) \arrow[r, "G(f)"'] & G(Y)
    \end{tikzcd}
\]

% =====================================================
\section{Tensor Notation}
\label{sec:tensor-notation}
% =====================================================

\subsection{Einstein Summation}

Repeated indices imply summation. The squared norm and metric operations:
\[
    \|\mathbf{v}\|^2 = v^i v_i = g_{ij} v^i v^j, \qquad
    v_i = g_{ij} v^j, \qquad v^i = g^{ij} v_j
\]

The Riemann curvature tensor symmetry:
\[
    R_{ijkl} = -R_{jikl} = -R_{ijlk} = R_{klij}
\]

\subsection{Covariant Derivative}

\[
    \nabla_m T^i{}_j = \partial_m T^i{}_j + \Gamma^i_{mk} T^k{}_j - \Gamma^k_{mj} T^i{}_k
\]

Christoffel symbols of the Levi-Civita connection:
\[
    \Gamma^k_{ij} = \frac{1}{2} g^{kl}\left(\partial_i g_{jl} + \partial_j g_{il} - \partial_l g_{ij}\right)
\]

\subsection{Bianchi Identities}

First (torsion-free): $R^i{}_{jkl} + R^i{}_{lkj} + R^i{}_{kjl} = 0$. Second:
\[
    \nabla_m R^i{}_{jkl} + \nabla_k R^i{}_{jlm} + \nabla_l R^i{}_{jmk} = 0
\]

% =====================================================
\section{Category Theory}
\label{sec:category-theory}
% =====================================================

A (covariant) functor $F\colon \mathcal{C} \to \mathcal{D}$ preserves composition and
identities:
\begin{align}
    F(g \circ f) &= F(g) \circ F(f), \label{eq:functor-comp}\\
    F(\Id_X) &= \Id_{F(X)}. \label{eq:functor-id}
\end{align}

An adjunction between $F\colon \mathcal{C} \to \mathcal{D}$ and
$G\colon \mathcal{D} \to \mathcal{C}$:
\[
    \Hom_{\mathcal{D}}(F(X), Y) \cong \Hom_{\mathcal{C}}(X, G(Y))
\]

The coequaliser of $f, g\colon A \rightrightarrows B$:
\[
    \begin{tikzcd}
        A \arrow[r, shift left, "f"] \arrow[r, shift right, "g"'] &
        B \arrow[r, "q"] & Q
    \end{tikzcd}
\]

The Yoneda lemma: $\Nat(\Hom_{\mathcal{C}}(-, X), F) \cong F(X)$ for
$F\colon \mathcal{C}^{\mathrm{op}} \to \mathbf{Set}$.

% =====================================================
\section{Proof Environments}
\label{sec:proof-environments}
% =====================================================

\begin{definition}[Metric Space]\label{def:metric}
A \emph{metric space} is a pair $(X, d)$ where $d\colon X \times X \to \mathbb{R}_{\geq 0}$ satisfies
non-negativity, symmetry, and the triangle inequality.
\end{definition}

\begin{theorem}[Bolzano--Weierstrass]\label{thm:bolzano-weierstrass}
Every bounded sequence in $\mathbb{R}^n$ has a convergent subsequence.
\end{theorem}

\begin{proof}[Proof of \autoref{thm:bolzano-weierstrass}]
By the bisection argument on a bounded interval, we extract a convergent subsequence.
\end{proof}

\begin{lemma}[Arzel\`{a}--Ascoli]\label{lem:arzela-ascoli}
$\mathcal{F} \subset C(X)$ is relatively compact iff uniformly bounded and equicontinuous.
\end{lemma}

\begin{proposition}\label{prop:lipschitz}
If $f\colon \mathbb{R} \to \mathbb{R}$ is Lipschitz with constant $L$, then $f$ is
absolutely continuous.
\end{proposition}

\begin{corollary}\label{cor:lipschitz}
Every Lipschitz function on $[a,b]$ is differentiable almost everywhere.
\end{corollary}

\begin{example}\label{ex:euclidean}
$d(x,y) = \sqrt{\sum_{i=1}^n (x_i - y_i)^2}$ defines a metric on $\mathbb{R}^n$.
\end{example}

\begin{remark}
Bolzano--Weierstrass fails in infinite-dimensional spaces; use weak compactness.
\end{remark}

% =====================================================
\section{Multi-line Aligned Equations}
\label{sec:aligned-equations}
% =====================================================

\subsection{The \texttt{align} Environment}

Derivation of the quadratic formula:
\begin{align}
    ax^2 + bx + c &= 0, \label{eq:quad-start}\\
    \left(x + \frac{b}{2a}\right)^2 &= \frac{b^2 - 4ac}{4a^2}, \label{eq:quad-square}\\
    x &= \frac{-b \pm \sqrt{b^2 - 4ac}}{2a}. \label{eq:quad-end}
\end{align}

\subsection{The \texttt{alignat} Environment}

Maxwell's equations with explicit column alignment:
\begin{alignat}{3}
    &\text{(i)} \quad &&\nabla \cdot \mathbf{E} = \frac{\rho}{\varepsilon_0} \quad &&\text{(Gauss's law)} \label{eq:maxwell1}\\
    &\text{(ii)} \quad &&\nabla \cdot \mathbf{B} = 0 \quad &&\text{(No monopoles)} \label{eq:maxwell2}\\
    &\text{(iii)} \quad &&\nabla \times \mathbf{E} = -\frac{\partial \mathbf{B}}{\partial t} \quad &&\text{(Faraday)} \label{eq:maxwell3}\\
    &\text{(iv)} \quad &&\nabla \times \mathbf{B} = \mu_0\mathbf{J} + \mu_0\varepsilon_0\frac{\partial \mathbf{E}}{\partial t} \quad &&\text{(Amp\`ere--Maxwell)} \label{eq:maxwell4}
\end{alignat}

\subsection{The \texttt{gathered} Environment}

\begin{equation}\label{eq:gathered-example}
    \begin{gathered}
        u_t = \kappa(u_{xx} + u_{yy}), \\
        u(x, y, 0) = f(x, y), \\
        u(x, y, t)\big|_{\partial\Omega} = 0.
    \end{gathered}
\end{equation}

% =====================================================
\section{Systems of Equations}
\label{sec:systems}
% =====================================================

\subsection{The \texttt{cases} Environment}

\begin{equation}\label{eq:heaviside}
    H(x) = \begin{cases}
        0 & \text{if } x < 0, \\
        \frac{1}{2} & \text{if } x = 0, \\
        1 & \text{if } x > 0.
    \end{cases}
\end{equation}

\subsection{The \texttt{dcases} Environment}

Double-sized fractions in piecewise definitions:
\begin{equation}\label{eq:dcases}
    f(x) = \begin{dcases}
        \frac{1}{2\sigma}\exp\!\left(-\frac{(x-\mu)^2}{2\sigma^2}\right)
            & \text{if } x \in [\mu - 3\sigma, \mu + 3\sigma], \\[6pt]
        0 & \text{otherwise}.
    \end{dcases}
\end{equation}

\subsection{The \texttt{empheq} Package}

Numbered boxed system for the Lorenz equations:
\begin{empheq}[box=\fbox]{align}
    \dot{x} &= \sigma(y - x) \label{eq:lorenz1}\\
    \dot{y} &= x(\rho - z) - y \label{eq:lorenz2}\\
    \dot{z} &= xy - \beta z \label{eq:lorenz3}
\end{empheq}

% =====================================================
\section{Matrix Notation}
\label{sec:matrices}
% =====================================================

\subsection{Standard Environments}

\[
    \mathbf{A} = \begin{pmatrix}
        a_{11} & a_{12} & a_{13} \\
        a_{21} & a_{22} & a_{23} \\
        a_{31} & a_{32} & a_{33}
    \end{pmatrix}, \qquad
    \begin{bmatrix}
        1 & 0 \\ 0 & 1
    \end{bmatrix}, \qquad
    \begin{vmatrix}
        a & b \\ c & d
    \end{vmatrix} = ad - bc
\]

\subsection{Small Matrices Inline}

The identity is $\bigl(\begin{smallmatrix} 1 & 0 \\ 0 & 1 \end{smallmatrix}\bigr)$.

\subsection{Augmented Matrices}

Row reduction:
\[
    \left(\begin{array}{ccc|c}
        2 & 1 & -1 & 8 \\
        -3 & -1 & 2 & -11 \\
        -2 & 1 & 2 & -3
    \end{array}\right)
    \xrightarrow{R_2 \leftarrow R_2 + \frac{3}{2}R_1}
    \left(\begin{array}{ccc|c}
        2 & 1 & -1 & 8 \\
        0 & \frac{1}{2} & \frac{1}{2} & 1 \\
        -2 & 1 & 2 & -3
    \end{array}\right)
\]

\subsection{Block Matrices}

\[
    \mathbf{M} = \begin{pmatrix}
        \mathbf{A} & \mathbf{B} \\
        \mathbf{C} & \mathbf{D}
    \end{pmatrix}, \qquad
    \mathbf{S} = \mathbf{D} - \mathbf{C}\mathbf{A}^{-1}\mathbf{B}
\]

% =====================================================
\section{Operator Definitions}
\label{sec:operators}
% =====================================================

Custom operators declared with \verb|\DeclareMathOperator|:
\begin{align}
    \tr(\mathbf{A}) &= \sum_{i=1}^n \lambda_i, \label{eq:trace}\\
    \det(\mathbf{A}) &= \prod_{i=1}^n \lambda_i, \label{eq:determinant}\\
    \rank(\mathbf{A}) &= \dim(\im(\mathbf{A})), \label{eq:rank}\\
    \ker(\mathbf{A}) &= \{\mathbf{x} : \mathbf{A}\mathbf{x} = \mathbf{0}\}. \label{eq:kernel}
\end{align}

The adjugate satisfies $\mathbf{A} \cdot \adj(\mathbf{A}) = \det(\mathbf{A}) \cdot \mathbf{I}_n$.

The error function: $\erf(x) = \frac{2}{\sqrt{\pi}} \int_0^x e^{-t^2}\,\mathup{d}t$.

% =====================================================
\section{Dimensional Analysis}
\label{sec:dimensional-analysis}
% =====================================================

Physical laws with proper units:
\begin{align}
    F &= m \cdot a = \qty{10}{\kilo\gram} \cdot \qty{9.81}{\meter\per\second\squared}
        = \qty{98.1}{\newton}, \label{eq:newton}\\
    E &= m c^2 = \qty{9.109e-31}{\kilo\gram} \cdot (\qty{2.998e8}{\meter\per\second})^2
        \approx \qty{8.187e-14}{\joule}, \label{eq:einstein}\\
    P &= \frac{F}{A} = \frac{\qty{100}{\newton}}{\qty{0.5}{\meter\squared}}
        = \qty{200}{\pascal}. \label{eq:pressure}
\end{align}

Measurement with uncertainty:
\[
    g = \qty{9.81}{\meter\per\second\squared} \pm \qty{0.02}{\meter\per\second\squared}, \qquad
    T = 2\pi\sqrt{\frac{L}{g}} \approx \qty{2.006}{\second}
\]

% =====================================================
\section{Chemical Equations}
\label{sec:chemical-equations}
% =====================================================

Combustion of methane:
\[
    \ch{CH4 + 2 O2 -> CO2 + 2 H2O}
\]

The Haber process:
\[
    \ch{N2 + 3 H2 <=>[$K_\text{eq}$] 2 NH3}
\]

Photosynthesis with state symbols:
\[
    \ch{6 CO2(g) + 6 H2O(l) ->[\text{light}] C6H12O6(aq) + 6 O2(g)}
\]

Net ionic equation for silver chloride precipitation:
\[
    \ch{Ag+(aq) + Cl-(aq) -> AgCl(s)}
\]

% =====================================================
\section{Logic Notation}
\label{sec:logic}
% =====================================================

Quantifiers and connectives:
\begin{align}
    \forall\, x \in \mathbb{R},\quad x^2 &\geq 0, \label{eq:universal}\\
    \exists\, x \in \mathbb{R},\quad x^2 &= 2, \label{eq:existential}\\
    \forall\, \varepsilon > 0,\;\exists\, \delta > 0 : |x - a| < \delta
        &\implies |f(x) - f(a)| < \varepsilon. \label{eq:epsilon-delta}
\end{align}

Logical equivalences:
\begin{align}
    P \iff Q &\equiv (P \implies Q) \land (Q \implies P), \label{eq:biconditional}\\
    \lnot(P \land Q) &\equiv \lnot P \lor \lnot Q, \label{eq:demorgan}\\
    \lnot(P \lor Q) &\equiv \lnot P \land \lnot Q. \label{eq:demorgan2}
\end{align}

% =====================================================
\section{Set Theory}
\label{sec:set-theory}
% =====================================================

Basic operations:
\begin{align}
    A \cup B &= \{x : x \in A \lor x \in B\}, \label{eq:union}\\
    A \cap B &= \{x : x \in A \land x \in B\}, \label{eq:intersection}\\
    A \setminus B &= \{x : x \in A \land x \notin B\}. \label{eq:difference}
\end{align}

The power set: $\mathcal{P}(S) = \{A : A \subseteq S\}$, $|\mathcal{P}(S)| = 2^{|S|}$.

De Morgan's laws for sets:
\[
    (A \cup B)^c = A^c \cap B^c, \qquad (A \cap B)^c = A^c \cup B^c.
\]

Cantor's theorem: $|S| < |\mathcal{P}(S)|$ for any set $S$.

% =====================================================
\section{Probability}
\label{sec:probability}
% =====================================================

Expectation and variance:
\begin{align}
    \mathbb{E}[X] &= \sum_{x} x\,\mathbb{P}(X = x) = \int_{-\infty}^{\infty} x f(x)\,\mathup{d}x, \label{eq:expectation}\\
    \Var(X) &= \mathbb{E}[(X - \mathbb{E}[X])^2] = \mathbb{E}[X^2] - (\mathbb{E}[X])^2. \label{eq:variance}
\end{align}

Bayes' theorem:
\[
    \mathbb{P}(A \mid B) = \frac{\mathbb{P}(B \mid A)\,\mathbb{P}(A)}{\mathbb{P}(B)}
\]

Common distributions:
\begin{align}
    X \sim \mathcal{N}(\mu, \sigma^2) &\implies f(x) = \frac{1}{\sigma\sqrt{2\pi}}
        \exp\!\left(-\frac{(x-\mu)^2}{2\sigma^2}\right), \label{eq:normal}\\
    X \sim \operatorname{Poisson}(\lambda) &\implies \mathbb{P}(X = k) = \frac{\lambda^k e^{-\lambda}}{k!}. \label{eq:poisson}
\end{align}

% =====================================================
\section{Statistics}
\label{sec:statistics}
% =====================================================

The $z$-test statistic:
\[
    z = \frac{\bar{X} - \mu_0}{\sigma / \sqrt{n}}
\]

A $95\%$ confidence interval:
\[
    \bar{X} \pm z_{\alpha/2} \cdot \frac{\sigma}{\sqrt{n}}
\]

Ordinary least-squares estimator:
\[
    \hat{\boldsymbol{\beta}} = (\mathbf{X}^{\top}\mathbf{X})^{-1}\mathbf{X}^{\top}\mathbf{y}
\]

Coefficient of determination:
\[
    R^2 = 1 - \frac{\sum_{i}(y_i - \hat{y}_i)^2}{\sum_{i}(y_i - \bar{y})^2}
\]

Maximum likelihood: $\ell(\theta) = \sum_{i=1}^n \log f(x_i \mid \theta)$, solved via
$\frac{\partial \ell}{\partial \theta} = 0$.

% =====================================================
\section{Optimization}
\label{sec:optimization}
% =====================================================

The Lagrangian for equality constraints:
\[
    \mathcal{L}(\mathbf{x}, \boldsymbol{\lambda}) = f(\mathbf{x}) + \sum_{i=1}^m \lambda_i\, g_i(\mathbf{x})
\]

KKT conditions for $\min f(\mathbf{x})$ s.t.\ $g_i(\mathbf{x}) \leq 0$, $h_j(\mathbf{x}) = 0$:
\begin{align}
    \nabla f(\mathbf{x}^*) + \sum_{i} \mu_i \nabla g_i(\mathbf{x}^*) + \sum_{j} \lambda_j \nabla h_j(\mathbf{x}^*) &= \mathbf{0}, \label{eq:kkt1}\\
    \mu_i\, g_i(\mathbf{x}^*) &= 0, \quad \mu_i \geq 0. \label{eq:kkt2}
\end{align}

Convexity: $f(\theta\mathbf{x} + (1-\theta)\mathbf{y}) \leq \theta f(\mathbf{x}) + (1-\theta) f(\mathbf{y})$.

Minimax: $\min_{\mathbf{x}} \max_{\mathbf{y}} \; \mathbf{x}^{\top} \mathbf{A} \mathbf{y}$.

% =====================================================
\section{Control Theory}
\label{sec:control-theory}
% =====================================================

State-space representation of an LTI system:
\begin{align}
    \dot{\mathbf{x}}(t) &= \mathbf{A}\mathbf{x}(t) + \mathbf{B}\mathbf{u}(t), \label{eq:state-space}\\
    \mathbf{y}(t) &= \mathbf{C}\mathbf{x}(t) + \mathbf{D}\mathbf{u}(t). \label{eq:output}
\end{align}

Transfer function: $G(s) = \mathbf{C}(s\mathbf{I} - \mathbf{A})^{-1}\mathbf{B} + \mathbf{D}$.

Second-order system: $G(s) = \frac{\omega_n^2}{s^2 + 2\zeta\omega_n s + \omega_n^2}$.

PID controller:
\[
    u(t) = K_p\, e(t) + K_i \int_0^t e(\tau)\,\mathup{d}\tau + K_d \frac{\mathup{d}e(t)}{\mathup{d}t}
\]

% =====================================================
\section{Fourier Analysis}
\label{sec:fourier}
% =====================================================

The Fourier transform and its inverse:
\begin{align}
    \hat{f}(\xi) &= \int_{-\infty}^{\infty} f(x)\, e^{-2\pi i \xi x}\,\mathup{d}x, \label{eq:fourier}\\
    f(x) &= \int_{-\infty}^{\infty} \hat{f}(\xi)\, e^{2\pi i \xi x}\,\mathup{d}\xi. \label{eq:inverse-fourier}
\end{align}

Fourier series for a periodic function with period $T$:
\[
    f(t) = \frac{a_0}{2} + \sum_{n=1}^{\infty} \left(
        a_n \cos\!\left(\frac{2\pi n t}{T}\right) + b_n \sin\!\left(\frac{2\pi n t}{T}\right)
    \right)
\]

Parseval's theorem: $\int |f(x)|^2\,\mathup{d}x = \int |\hat{f}(\xi)|^2\,\mathup{d}\xi$.

Convolution theorem: $\widehat{f * g}(\xi) = \hat{f}(\xi) \cdot \hat{g}(\xi)$.

% =====================================================
\section{Differential Geometry}
\label{sec:differential-geometry}
% =====================================================

A smooth $n$-manifold $\mathcal{M}$ has an atlas of charts
$\{(U_\alpha, \varphi_\alpha)\}$ with $\varphi_\alpha\colon U_\alpha \to \mathbb{R}^n$.

Curvature tensor of a connection:
\[
    R(X, Y) = \nabla_X \nabla_Y - \nabla_Y \nabla_X - \nabla_{[X,Y]}
\]

Geodesic equation: $\ddot{x}^k + \Gamma^k_{ij}\, \dot{x}^i \dot{x}^j = 0$.

Stokes' theorem for an $(n{-}1)$-form $\omega$ on a compact oriented $n$-manifold:
\[
    \int_{\mathcal{M}} \mathup{d}\omega = \int_{\partial\mathcal{M}} \omega
\]

% =====================================================
\section{Algebraic Structures}
\label{sec:algebra}
% =====================================================

A group $(G, \cdot)$ satisfies closure, associativity, identity, and invertibility.
The symmetric group $S_n$ has $|S_n| = n!$.

\begin{theorem}[Lagrange]\label{thm:lagrange}
If $H \leq G$ is a subgroup of a finite group, then $|H|$ divides $|G|$.
\end{theorem}

A ring $(R, +, \cdot)$ has $(R, +)$ abelian and $(R, \cdot)$ a monoid distributing over addition.
A field is a commutative ring with multiplicative inverses for all nonzero elements.

Module axioms for an $R$-module $M$:
\begin{align}
    r(m + n) &= rm + rn, \label{eq:module1}\\
    (r + s)m &= rm + sm, \label{eq:module2}\\
    (rs)m &= r(sm), \label{eq:module3}\\
    1 \cdot m &= m. \label{eq:module4}
\end{align}

First isomorphism theorem: $G / \ker(\varphi) \cong \im(\varphi)$ for a homomorphism $\varphi$.

% =====================================================
\section{Number Theory}
\label{sec:number-theory}
% =====================================================

Fundamental theorem of arithmetic: every $n > 1$ factors uniquely as
$n = p_1^{a_1} p_2^{a_2} \cdots p_k^{a_k}$. Prime number theorem: $\pi(x) \sim \frac{x}{\ln x}$.

\begin{definition}[Congruence]\label{def:congruence}
$a \equiv b \pmod{n}$ iff $n \mid (a - b)$.
\end{definition}

\begin{theorem}[Fermat]\label{thm:fermat}
If $p$ is prime and $\gcd(a, p) = 1$, then $a^{p-1} \equiv 1 \pmod{p}$.
\end{theorem}

\begin{theorem}[Euler]\label{thm:euler}
If $\gcd(a, n) = 1$, then $a^{\varphi(n)} \equiv 1 \pmod{n}$.
\end{theorem}

The Legendre symbol:
\[
    \left(\frac{a}{p}\right) = \begin{cases}
        1 & \text{if } a \text{ is a quadratic residue mod } p, \\
        -1 & \text{if } a \text{ is a quadratic non-residue mod } p, \\
        0 & \text{if } p \mid a.
    \end{cases}
\]

\begin{example}
$x^5 - x - 1$ is irreducible over $\mathbb{Q}$ with $\Gal \cong S_5$ (not solvable),
so the general quintic has no radical solution.
\end{example}

% =====================================================
\section{Conclusion}
\label{sec:conclusion}
% =====================================================

This gallery demonstrates 20 advanced mathematical typesetting techniques, from
commutative diagrams and tensor calculus to Galois theory and Fourier analysis.
The combination of \texttt{mathtools}, \texttt{empheq}, \texttt{tikz-cd},
\texttt{siunitx}, and \texttt{chemmacros} provides a comprehensive toolkit
for mathematical and scientific writing.

\end{document}
